\section{从 Gauss 曲率到统计与数据流形}
\label{sec:12-Real-Analysis-Support/11-Curve-and-Surface-Geometry/03}

包一个圆筒形的礼盒，纸能服服帖帖贴上去，一道褶都不用打。包一个皮球，无论怎么摆弄，纸上必定出现褶子或者要剪开一块。纸的性质是不肯拉伸；圆筒不逼它拉伸，球逼它。

两个形状在 \(\R^{3}\) 里都弯着，包装纸却能分辨它们。这说明它们的差别不在``看上去弯不弯''，而在某种纸自己就能感知的性质上。上一节留下的悬案正是这件事：Gauss 曲率的定义里用到了第二基本形式，按说是外在的，可它偏偏能被一张不肯拉伸的纸测出来。

\begin{center}
\begin{tikzpicture}[>=stealth,scale=0.72]
% 左：圆柱面，可展
\draw[gray!60] (0,0.28) arc (180:360:0.8 and 0.28);
\draw[gray!35,densely dashed] (0,0.28) arc (180:0:0.8 and 0.28);
\draw[gray!60] (0,2.5) ellipse (0.8 and 0.28);
\draw[gray!60] (0,0.28) -- (0,2.5);
\draw[gray!60] (1.6,0.28) -- (1.6,2.5);
\foreach \a in {0.45,0.8,1.15} \draw[gray!30] (\a,0.13) -- (\a,2.35);
\node[font=\small] at (0.8,-0.55) {圆柱面，\(K=0\)};
\draw[->,thick,gray!70!black] (2.1,1.35) -- (3.1,1.35);
\node[font=\small,gray!70!black] at (2.6,1.75) {剪开摊平};
% 摊平后的矩形
\draw[orange!85!black] (3.5,0.28) rectangle (5.6,2.5);
\foreach \a in {4.0,4.55,5.1} \draw[orange!35] (\a,0.28) -- (\a,2.5);
\foreach \b in {1.02,1.76} \draw[orange!35] (3.5,\b) -- (5.6,\b);
\node[font=\small] at (4.55,-0.55) {不拉伸，不撕裂};
% 右：球面，不可展
\draw[gray!60] (8.6,1.4) circle (1.15);
\begin{scope}
  \clip (8.6,1.4) circle (1.15);
  \draw[gray!30] (8.6,1.4) ellipse (1.15 and 0.34);
  \draw[gray!30] (8.6,1.94) ellipse (1.0 and 0.28);
  \draw[gray!30] (8.6,0.86) ellipse (1.0 and 0.28);
  \draw[gray!30] (8.6,1.4) ellipse (0.40 and 1.15);
  \draw[gray!30] (8.6,1.4) ellipse (0.80 and 1.15);
\end{scope}
\node[font=\small] at (8.6,-0.55) {球面，\(K=1/R^{2}\)};
\draw[->,thick,gray!70!black] (10.0,1.35) -- (11.0,1.35);
\node[font=\small,gray!70!black] at (10.5,1.75) {试图摊平};
% 摊平后的瓣（必留缝隙）
\foreach \k in {0,1,2,3} {
  \draw[blue!70!black] plot[domain=0:1,samples=30]
    ({11.95+0.78*\k+0.30*sin(180*\x)},{0.28+2.22*\x});
  \draw[blue!70!black] plot[domain=0:1,samples=30]
    ({11.95+0.78*\k-0.30*sin(180*\x)},{0.28+2.22*\x});
}
\node[font=\small] at (13.15,-0.55) {缝隙合不拢};
\end{tikzpicture}
\end{center}

\emph{圆柱面剪开就能平铺，球面不行；差别不在外观，在一个内部就能测出的数。}

\subsection{绝妙定理：Gauss 曲率只依赖第一基本形式}

\begin{theorem}[Theorema Egregium]
曲面的 Gauss 曲率 \(K\) 可以只用第一基本形式的分量 \(E,F,G\) 及其一阶、二阶偏导数表达出来。特别地，局部等距的两张曲面在对应点处 \(K\) 相等。
\end{theorem}

在正交参数化（\(F=0\)）下，公式长这样：

\[
K=-\frac{1}{2\sqrt{EG}}\left[
\partial_{u}\!\left(\frac{G_{u}}{\sqrt{EG}}\right)
+\partial_{v}\!\left(\frac{E_{v}}{\sqrt{EG}}\right)
\right].
\]

一般参数化下的表达式更长，但同样只出现 \(E,F,G\) 和它们的导数，第二基本形式的分量一个也没有。

证明骨架分三步。先把 \(X_{uu},X_{uv},X_{vv}\) 按基 \(X_{u},X_{v},n\) 展开，切向分量的系数是 Christoffel 记号，它们只由 \(E,F,G\) 及其一阶导决定，法向分量的系数则是 \(e,f,g\)。接着利用混合偏导可交换 \(\partial_{v}X_{uu}=\partial_{u}X_{uv}\)，把两边按同一组基展开并逐分量比较。切向分量给出的等式（Gauss 方程）恰好把 \(eg-f^{2}\) 表达成 Christoffel 记号的导数与乘积的组合，除以 \(EG-F^{2}\) 即得 \(K\)。条件对账：只用到参数化二阶连续可微以便交换偏导，以及 \(X_{u},X_{v},n\) 处处构成基，也就是正则性。没有任何一步依赖曲面在 \(\R^{3}\) 中的具体位置。

\textbf{一个看起来关于``怎样嵌入外部空间''的量，其实纯粹内蕴。} 这就是 Gauss 把它叫做``绝妙''的理由：\(e,f,g\) 各自都会随嵌入方式改变，组合 \(eg-f^{2}\) 除以 \(EG-F^{2}\) 之后，所有对外部的依赖恰好抵消干净。

\subsection{住在曲面上的人不用离开曲面就能测出它}

内蕴意味着可操作。设想一只蚂蚁被困在曲面上，没有第三维的概念，手里只有一根尺。它照样能测出 \(K\)。

一个办法是量圆周。以 \(p\) 为中心、沿曲面量出半径为 \(r\) 的测地圆，它的周长满足

\[
C(r)=2\pi r\left(1-\frac{K(p)}{6}r^{2}+O(r^{4})\right).
\]

平面上 \(C(r)=2\pi r\) 分毫不差。正曲率处周长偏小——球面上以北极为心、测地半径 \(r\) 的圆，其真实周长是 \(2\pi R\sin(r/R)\)，确实比 \(2\pi r\) 小。负曲率处周长偏大。蚂蚁只要把测得的周长和 \(2\pi r\) 比一比，就能读出 \(K\) 的符号，量得够准还能读出大小。

另一个办法是量三角形。以测地线为边的三角形，内角和减去 \(\pi\) 等于 \(K\) 在三角形内的积分。球面上的三角形内角和大于 \(\pi\)，马鞍面上小于 \(\pi\)。这两件事都不需要抬头看第三维。

回到那张包装纸。纸的内蕴几何是平的，\(K\equiv0\)；圆柱面的 \(K\) 也处处为零，两者局部等距，所以贴合毫不费力。球面的 \(K=1/R^{2}\neq0\)，与纸不等距，绝妙定理直接判定：不撕不拉是贴不上去的。褶子不是手笨，是定理。

\subsection{地图必定失真，这是定理不是技术问题}

同一个论证解释了一件早先出现过的事。\hyperref[sec:12-Real-Analysis-Support/08-Manifold-Basics/01]{``流形是局部像欧氏空间的集合''}说球面画不成一张不失真的平面地图，当时的理由是拓扑的——紧集不同胚于开集，所以一张卡盖不住整个球面。绝妙定理给出的是更锋利的版本：哪怕只画一小块，只要要求长度全部忠实，也做不到。

推理只有一行。若某张地图在一小片区域上保持所有长度，那它就是一个局部等距，于是两边的 \(K\) 必须相等。可是平面上 \(K=0\)，球面上 \(K=1/R^{2}>0\)。矛盾。

所以任何地图投影都必须放弃点什么。墨卡托保角度放弃面积，等积投影保面积放弃角度，等距圆柱投影只在特定方向上保长度。选投影，就是选牺牲哪一项——这不是制图学的懒惰，是曲率不为零逼出来的取舍。

顺带把球面的账目算全：取 \(R=1\)，第一基本形式是 \(\diag(1,\sin^{2}\theta)\)，第二基本形式是 \(-\mathrm{I}\)，形状算子是恒等映射的相反数，两个主曲率都是 \(-1\)，于是 \(K=1\)、\(H=-1\)。法向取反会让 \(H\) 变号而 \(K\) 不变，这也是 \(K\) 更稳健的一个侧面。测地线是大圆，两条经线在赤道处平行、到极点却相交——正曲率让起初平行的测地线汇聚，负曲率让它们发散，零曲率让它们保持平行。这是曲率最直观的可观测效应。

\subsection{落到数据上：曲率控制局部线性近似的误差}

现在说数据，而且要说得克制。

有一件事是定理。切空间是流形的一阶近似，沿切方向 \(w\) 走出距离 \(t\) 时，曲面偏离切平面的量首项是 \(\tfrac12\,\mathrm{II}(w,w)t^{2}\)。局部线性近似的误差由二阶量控制，量级是 \(t^{2}/\rho\)，其中 \(\rho\) 是曲率半径。这句话在任何光滑子流形上都成立。

由它可以推出几条量级估计，注意它们是估计而非定理。局部主成分分析、局部线性嵌入取的邻域，半径应当明显小于曲率半径，否则那个``局部线性块''装不下真实的形状；曲率大的区域需要更小的邻域，而更小的邻域需要更密的采样，两个要求在样本有限时会打架。Isomap 用邻接图的最短路近似测地距离，弦与弧的差距同样由曲率控制，弯得越急的地方图最短路偏得越多。内蕴弯曲的数据流形无法被等距地铺进同维数的欧氏空间，所以任何保距降维都必然有失真——这与地图那件事是同一个定理的不同穿法。

再往前的说法就属于启发了。``某个数据集的流形曲率大''这种断言，在实践中很难验证：曲率是二阶量，估计它所需的样本量比估计切空间高一个数量级，而且对噪声极其敏感——噪声尺度一旦接近曲率半径，估出来的曲率基本是噪声的曲率。样本不充裕时，先假定低曲率、把邻域取小，比强行估计更稳妥。曲率在这里的价值主要是提供一套语言，让``邻域该取多大''``降维为什么失真''这类问题有地方安放，而不是提供一个可以直接算出来的数字。

统计流形那边的情形类似，而且多一层讲究。Fisher 度量把参数族变成黎曼流形之后，曲率才有定义；但曲率还依赖联络的选择，Levi-Civita 联络与信息几何中的 \(\alpha\)-联络给出不同的答案，统计含义也不同。指数族在对偶联络下是平坦的，这正是自然参数与期望参数两套``平直坐标''存在的几何原因；混合模型则不平坦。说一句``这个统计流形是弯的''，不指明在哪个联络下，等于没说。\hyperref[ch:105]{``信息几何：参数空间中的分布几何''}处理的就是这一层。

\subsection{本单元到此为止}

从曲线的一个数，到曲面的两个基本形式，再到绝妙定理，这条线索反复在讲同一件事：把依赖坐标、依赖嵌入的东西剥掉之后，剩下的才是几何。曲率是这条线索的顶点——它最初被写成一个明显依赖外部空间的表达式，最终被证明属于曲面自己。

对本书而言，需要的几何到此为止。联络、Jacobi 场、示性类属于黎曼几何本身，那是研究弯曲这件事的学问；这里只是借它的工具，把``这个流形是弯的''从一句比喻变成一个可以计算、可以证伪的句子。有些性质无法通过换一种嵌入方式消除——记住这一条，比记住任何一个曲率公式都有用。
